\section{Completed single-primary boundary calculation and the continuing mixed-support calculation}
\label{sec:bc-single-primary-update}

The full single-primary calculation SP.1--36 and SPF.1--29,
included in this edition as
\emph{The actual single-primary marked family and its complete boundary
monodromy} and its finite-field companion, sharpens the preceding
connection and determinant calculations on their exact primary
subcase.  The formulas in this section are reused from that complete
proof, with the original source, constant, unit and basis retained.
They replace the earlier statement that this subcase's local inertia
and periods still required calculation.

\subsection{The exact primary specialization, connection and period matrix}

Choose an actual nontrivial zero \(\rho\) of \(g=2\xi\), with its
complete order \(m\geq1\).  In this subsection fix \(t=0\) and put
\[
 n=m+1,\qquad h(s)=(s-\rho)^m,\qquad
 y=s-\rho,\qquad
 \Phi_h(s)=\frac{(s-\rho)^n-(-\rho)^n}{n}
          =\frac{y^n}{n}+c,\qquad
 c=-\frac{(-\rho)^n}{n}.
\]
Thus \(\Phi_h(0)=0\) and \(\Phi_h'=h\); the constant \(c\) has its
original value.  In the ordered rows
\(e_s=(1,s,\ldots,s^{m-1})\) and
\(e_y=(1,y,\ldots,y^{m-1})\), the exact coefficient map is
\[
 e_s=e_yP,\qquad
 P_{ab}=\binom ba\rho^{b-a}\quad(a\leq b),\qquad
 (P^{-1})_{ab}=\binom ba(-\rho)^{b-a}\quad(a\leq b),
\]
with entries zero when \(a>b\).  Substitution in the binomial theorem
proves both maps and their inverse relation; \(\det P=1\).

The original complex has differential
\(L_u=u\partial_s+(s-\rho)^m\).
Its leading-degree term gives the free degree-one remainder module
\(\bigoplus_{b=0}^{m-1}\mathbb C[u][s^bds]\), whose marked fibre is
\(\mathbb C[s]/h\,ds\), with the full primary order.  In these unchanged
coefficient columns the exact division and its derivative are
\[
 \Phi_h(s)s^b
   =h(s)\frac{(s-\rho)s^b}{n}+cs^b,\qquad
 \left(\frac{(s-\rho)s^b}{n}\right)'
   =\frac{(b+1)s^b-b\rho s^{b-1}}n .
\]
Subtracting \(L_u((s-\rho)s^b/n)\) proves
\[
 C_s=cI_m,\qquad
 (B_s)_{bb}=\frac{b+1}{n},\qquad
 (B_s)_{b-1,b}=-\frac{b\rho}{n},\qquad
 \nabla_u=\partial_u-\frac{cI_m}{u^2}+\frac{B_s}{u}.
\]
All other entries of \(B_s\) vanish.  The degree-zero lift is still
\(\partial_u-\Phi_h/u^2+1/u\); it intertwines \(L_u\) with the
degree-one lift \(\partial_u-\Phi_h/u^2\).  Writing
\(\beta_a=(a+1)/n\) gives the exact conjugation
\[
 PB_sP^{-1}=\operatorname{diag}(\beta_0,\ldots,\beta_{m-1}),
 \qquad \operatorname{Tr}B_s=m/2.
\]
These are SP.1--9, including the original coefficient and chain maps.

On a fixed branch of \(\log u\), retain the ordered outward rays of
angles \((\arg u+\pi)/n+2\pi j/n\) and their differences
\(\Gamma_j=-\ell_0+\ell_j\), \(1\leq j\leq m\).
Set
\[
 \zeta=e^{2\pi i/n},\qquad
 W_{ja}=\zeta^{j(a+1)}-1,\qquad
 d_a=n^{\beta_a-1}e^{\pi i\beta_a}\Gamma(\beta_a),\qquad
 D(u)=\operatorname{diag}(d_a u^{\beta_a}).
\]
The complete original period matrix, with original monomial columns,
is
\[
 \Pi_s(u)=\left(\int_{\Gamma_j}e^{\Phi_h(s)/u}s^b\,ds\right)_{jb}
         =e^{c/u}W D(u)P .
\]
SP.10--14 prove this equality on the original contours: translating
their junction by \(\rho\) leaves connector integrals tending to zero
by
\(\operatorname{Re}(\Phi_h/u)
=-R^n/(n|u|)+O(R^{n-1})+O(1)\); the finite junction paths cancel with
their original orientations.  The radial substitution then gives
exactly the \(d_a\) above.  The matrix \(W\) is invertible by the
degree-\((n-1)\) polynomial root argument in SP.14.  In particular,
\[
 \det\Pi_s
 =K_m u^{m/2}e^{mc/u},\qquad
 K_m=n^{-m/2}e^{\pi im/2}
       \left(\prod_{a=0}^{m-1}\Gamma(\beta_a)\right)\det W\ne0 .
\]
Thus the preceding determinant has its complete matrix refinement
on this precise subcase, with all Gamma factors, contour phases,
Pascal coefficients and the original scalar retained.

\subsection{The computed boundary and its exact marked-source maps}

The transported cycles satisfy
\(\Gamma_j\mapsto\Gamma_{j+1}-\Gamma_1\) for \(j<m\) and
\(\Gamma_m\mapsto-\Gamma_1\).  Consequently SP.16--21 give
\[
 M_{\rm cyc}
 =W\operatorname{diag}(\zeta,\zeta^2,\ldots,\zeta^m)W^{-1},
 \qquad M_{\rm cyc}^n=I,\qquad
 \ker(M_{\rm cyc}-I)=0 .
\]
The horizontal coefficient convention is determined by the exact
fundamental matrix
\[
 H_s(u)=P^{-1}e^{-c/u}\operatorname{diag}(u^{-\beta_a}).
\]
It has inverse eigencharacters, with the same zero invariant space.
Under \(u=v^n\), the displayed connection becomes
\[
 \partial_v-\frac{nc}{v^{n+1}}I_m
     +\frac1v\operatorname{diag}(1,\ldots,m).
\]
The meromorphic change
\(x_y=\operatorname{diag}(v^{-1},\ldots,v^{-m})z\)
removes precisely these integer residues and retains the common
irregular factor \(e^{-c/v^n}\).  This exact fundamental solution
has identity monodromy after the cover; its unipotent logarithm is
\(N_{\rm boundary}=0\).  Its sector restrictions agree on overlaps
with the same branch, so there are no additional Stokes factors in
this one-critical-value calculation.  Both ordinary invariant
stalks and the derived local invariant and coinvariant groups vanish
before the cover, since \(M_{\rm cyc}-I\) is invertible.

On the same marked fibre let \(J_m\) be multiplication by \(y\), with
\((J_m)_{a+1,a}=1\).  The original multiplication and its precise
relation to the boundary connection are
\[
 A_s=P^{-1}(\rho I_m+J_m)P,\qquad
 N_{\rm ar}=P^{-1}J_mP,\qquad
 [B_s,A_s]=\frac1nN_{\rm ar},\qquad
 [\nabla_u,A_s]=\frac{N_{\rm ar}}{nu}.
\]
With \(D_0=\operatorname{diag}(d_a)\) and
\(K_0=WD_0J_mD_0^{-1}W^{-1}\), the complete period transport is
\[
 \Pi_sA_s\Pi_s^{-1}=\rho I_m+u^{1/n}K_0,\qquad
 M_{\rm cyc}K_0M_{\rm cyc}^{-1}=\zeta K_0 .
\]
The coefficient calculation on each monomial proves these
commutators and conjugations, including the last nilpotent column.
They retain the original arithmetic nilpotent through its exact
nonhorizontal map; see SP.22--24.

For the entire quotient \(v_h=g/h\), retain every coefficient
\[
 a_j=\frac{v_h^{(j)}(\rho)}{j!}
     =\frac{g^{(m+j)}(\rho)}{(m+j)!},\qquad 0\leq j<m,
 \qquad a_0\ne0 .
\]
Its marked multiplication and inverse are
\[
 U_s=P^{-1}\left(\sum_{j=0}^{m-1}a_jJ_m^j\right)P,\qquad
 U_s^{-1}=P^{-1}\left(\sum_{j=0}^{m-1}b_jJ_m^j\right)P,
\]
\[
 b_0=a_0^{-1},\qquad
 b_r=-a_0^{-1}\sum_{j=1}^r a_jb_{r-j}.
\]
Multiplication of these finite series proves their inverse relation,
using \(J_m^m=0\).  Its complete period and connection images are
\[
 \Pi_sU_s\Pi_s^{-1}
    =\sum_{j=0}^{m-1}a_j u^{j/n}K_0^j,\qquad
 [\nabla_u,U_s]
    =\frac1{nu}P^{-1}
          \left(\sum_{j=1}^{m-1}j a_jJ_m^j\right)P.
\]
These retain the full Taylor unit, rather than only its leading
coefficient.

The original Euler/theta square remains, with
\(D=-x\partial_x\) and \(\mathcal MF_h=v_h\),
\[
 \alpha_h(P)=P(D)\phi_*,\qquad
 \mathcal T_h(P)=P(D)F_h,\qquad
 \Theta\alpha_h(P)=\mathcal T_h(hP),\qquad
 J_h\mathcal T_h(P)=U_s[P]_s.
\]
Here
\(\phi_*=(4\pi^2x^4-6\pi x^2)e^{-\pi x^2}\),
\(\Theta\phi=2\sum_{r\geq1}\phi(rx)\), and
\(\mathcal M\Theta\phi_*=g\), with the original function spaces and
Mellin convention.  The degree-one map is the retained
\(\sigma_hU_s:E_h\hookrightarrow Q_{\rm ar}\).
The exact translated source vector is
\(\mathcal T_h(y^b)=(D-\rho)^bF_h\), and at general \(u\) the
source differential is the specified
\(\Theta+u\mathcal T_h\partial_s\alpha_h^{-1}\) on
\(\alpha_h(\mathbb C[s])\).
For \(v_{\rm pol}=\sum_{j<m}a_j(s-\rho)^j\), the original differential
gives the full representative defect
\[
 v_{\rm pol}L_uQ=L_u(v_{\rm pol}Q)-u v_{\rm pol}'Q.
\]
If \(v_{\rm pol}p=hq+r\) for an original remainder \(p\), its
differential remainder is \(r-uq'\), whereas the constant
coefficient extension \(U_s\) returns \(r\).
SP.25--31 prove these identities and retain precisely the derivative
term needed away from \(u=0\).

The exact finite-cover lattice map is also available:
\[
 e^{-c/v^n}\Pi_s(v^n)
    =WD_0\operatorname{diag}(v,v^2,\ldots,v^m)P .
\]
After the two displayed constant invertible changes, its cokernel is
\(\bigoplus_{a=0}^{m-1}\mathbb C\{v\}/(v^{a+1})\), of length
\(m(m+1)/2\).  Its image lattice has rank-\(m\) fibre modulo \(v\),
while its inclusion in the ambient constant lattice induces zero
modulo \(v\).  This is the complete lattice map between the retained
marked fibre and the computed local boundary; the factored
exponential remains an irregular factor.  See SP.32--33.

In the unchanged source metric of the preceding calculation, the
full matrix refinement is
\[
 T_{\rm per}(u)
 =G_\theta^{-1}\Pi_s(u)^*\Pi_s(u)
 =e^{2\operatorname{Re}(c/u)}
   G_\theta^{-1}P^*D(u)^*W^*WD(u)P .
\]
It is positive and self-adjoint in \(G_\theta\), with the complete
multiplier \(g/h\) in that Gram.  For \(N\geq m-1\), the enlarged
polynomial source has the exact period factorization and kernel
\[
 \operatorname{Per}_N=\Pi_sJ_{u,N},\qquad
 \ker J_{u,N}=L_u(\mathbb C[s]_{\leq N-m}).
\]
Thus its period Gram has rank \(m\); on the displayed remainder
frame it is positive definite.  SP.34--36 prove this by the original
Mellin source and the endpoint calculation for
\(u\,d(e^{\Phi_h/u}Q)\).  The original relation space and source norm
therefore remain in subsequent restriction calculations.

\subsection{The exact finite-field boundary on the same primary phase}

For the finite-field statement retain the specified coefficient
homomorphism \(R\to k_0=\mathbb F_{Q_0}\), with \(p>n\),
including \(\rho\), \(1/n!\), and every coefficient of the full
Taylor unit and its inverse.  Use the explicit constant extension
\[
 k=\mathbb F_Q,\qquad Q=Q_0^d,\qquad
 d=\operatorname{ord}_{(\mathbb Z/n\mathbb Z)^\times}(Q_0),
\]
so \(n\mid Q-1\).  Fix \(\ell\ne p\), geometric Frobenius, and the
nontrivial additive character \(\psi\) obtained from the specified
one on \(k_0\) by the field trace.  SPF.1--29 compute the actual sheaf
\(\mathcal F=R^1\pi_{s!}\mathcal L_\psi(\Phi_h(s)/u)\).
With
\[
 X_n=\{\chi:k^\times\to\overline{\mathbb Q}_\ell^\times:
             \chi^n=1\},\qquad
 G_k(\chi,\psi)=\sum_{z\in k^\times}\chi(z)\psi(z),
\]
and \(\mathcal A_\gamma\) denoting the geometrically constant line
with geometric Frobenius eigenvalue \(\gamma\), the exact
parameter-sheaf isomorphism is
\[
 \mathcal F\simeq
 \mathcal L_\psi(c/u)\otimes
 \bigoplus_{\substack{\chi\in X_n\\\chi\ne1}}
   \mathcal A_{-\chi(n)G_k(\chi,\psi)}
                  \otimes\mathcal K_\chi(u).
\]
The Kummer convention has trace \(\chi\).  SPF.6--19 prove the
isomorphism using the original translation, the finite power map
\(y\mapsto y^n\), its character projectors including the zero stalk,
and compact-support base change.  The minus sign is the degree-one
cohomological sign; \(\chi(n)\) is the original coefficient factor.
Both it and the common Artin--Schreier factor remain.

If the specialized \(c\ne0\), the common wild character has break
one, so
\[
 V^{I_0}=0,\qquad \operatorname{Swan}_0(\mathcal F)=m.
\]
If \(c=0\), the nontrivial tame Kummer characters instead give
\[
 V^{I_0}=0,\qquad \operatorname{Swan}_0(\mathcal F)=0.
\]
Writing \(\psi(x)=\vartheta(\operatorname{Tr}_{k/\mathbb F_p}(bx))\),
the explicit cover \(u=v^n\), followed when needed by
\(z^p-z=bc/v^n\), kills all geometric inertia.  Its restricted
unipotent logarithm is zero, while its restricted invariant space
is the full rank \(m\).  SPF.22--27 prove the conductor, original
invariant space and finite-cover statements by the exact local
ramification calculation.
Finally SPF.28--29 prove
\(|G_k(\chi,\psi)|^2=Q\) directly for every nontrivial \(\chi\);
all local eigenvalues consequently have the original weight one.
These are completed local and weight calculations for the specified
sheaf and coefficient specialization.

\subsection{Continuation through the full mixed-support family}

The preceding complete primary calculation now supplies actual
boundary matrices, period maps, conductors, unit defects and source
forms to the earlier family.  The original independent coefficient
faces, outer theta-leg labels, proper-source kernels and structural
specialization remain those already defined above; every constant
coefficient map retains the corresponding represented zero and its
label.

For several distinct primary critical values, the remaining
calculation includes the sector transitions and local extensions
between their computed blocks, and their exact maps through the
retained mixed faces.  On tensor powers the one-factor vanishing
cannot be propagated without the tensor character calculation:
the computed characters multiply under tensor product, and their
product can equal one.  The common finite-field additive factor
also becomes \(\mathcal L_\psi(rc/u)\) on the \(r\)-fold diagonal
tensor, retaining its actual residue characteristic.  These are
the precise character and coefficient maps to carry forward,
with the full original units and marked source maps.

The continuing induced relation-form calculation on the weighted
cyclic image retains the parameter defect
\eqref{eq:cyclic-defect}, the transverse map, its rectangular period
restriction, and all original source cross-pairings.  The explicit
primary formula for \(T_{\rm per}\) above can now enter that
calculation, while its complete coefficient Gram and restriction
maps remain.  An arithmetic Frobenius comparison and the uniform
four-volume estimate still require their own calculations on those
original objects; the completed primary boundary result has not
performed those additional steps.
